\documentclass[book.tex]{subfiles}
\begin{document}

\chapter{Linear Systems: Analytical Methods}
\label{chap:linear_analytical}
\noindent\rule{11cm}{0.4pt} \\
\textbf{Prerequisites:} \autoref{chap:linear_systems}, \autoref{chap:optimization} \\
\textbf{Difficulty Level:} * \\
\noindent\rule{11cm}{0.4pt}
\vspace{5mm}

In this chapter we will learn how to solve small sets of linear equations with graphical methods. We will then explain how to evaluate the determinant of a matrix and solve linear equations with Cramer's rule. Finally we will explain how to eliminate unknowns in equations by combining equations.
\section{Graphical Methods}
For a system of two variables, graphical method obtains solutions by plotting them on Cartesian coordinates with one axis corresponding to $x_1$ and the other to $x_2$. The values of $x_1$ and $x_2$ at the intersection of two lines represent the solution. Consider the example:
\begin{align}
3x_1 + 2x_2 &= 18\\
-x_1 + 2x_2 &= 2
\end{align}
Solving for $x_2$ in terms of $x_1$, we get
\begin{align}
x_2 = \dfrac{1}{2}x_1 + 1
\end{align}
\begin{marginfigure}
	\centering
	\includegraphics[width = 2.4in]{figures/part1b/linear_analytical/graphical.PNG}
	\centering
	\caption{Graphical solution of two simultaneous algebraic equations}
	\label{fig:1}
\end{marginfigure}
As shown in Figure \ref{fig:1}, two lines intersect at $x_1$ = 4 and $x_2$ = 3 which therefore is the solution for the example considered. 

For three simultaneous equations, each equation would be represented by a plane in a three-dimensional coordinate system. The point where the three planes intersect would represent the solution. For systems with more than three equations, graphical methods cannot be applied. Some other cases that can pose problems when solving sets of linear equations are as follows.

\begin{itemize}
	\item \textit{Singular systems}
	\begin{enumerate}
		\item When two equations represent parallel lines, there is no solution as they never intersect.
		\item When two lines are coincident, there is an infinite number of solutions.
	\end{enumerate}
	\item \textit{Ill-conditioned systems}: When the slopes are so close that it is difficult to detect the point of intersection. This system is very close to being singular. Also this system is extremely sensitive to round-off error.
\end{itemize}
\section{Determinants and Cramer's Rule}
Consider a system of three equations defined as $[A]$ $\left\{x\right\}$ = $\left\{b\right\}$ where $[A]$ is the coefficient matrix. The determinant of square matrix [A] is formed from its elements and is represented by D.
\begin{align}
[A] &=
\begin{bmatrix}
	a_{11} & a_{12} & a_{13} \\ 
	a_{21} & a_{22} & a_{23} \\ 
	a_{31} & a_{32} & a_{33}
\end{bmatrix}\\
D &= 
\begin{vmatrix}
a_{11} & a_{12} & a_{13} \\ 
a_{21} & a_{22} & a_{23} \\ 
a_{31} & a_{32} & a_{33}
\end{vmatrix}
\vspace{-6pt}
\end{align}
Determinant is a single number as compared to a matrix. For a system of two simultaneous equations, determinant is computed as 
\vspace{-8pt}
\begin{align}
D =
\begin{vmatrix}
a_{11} & a_{12} \\ 
a_{21} & a_{22}  
\end{vmatrix}
= a_{11} a_{22} - a_{12} a_{21}
\vspace{-8pt}
\end{align}
For the above matrix [A], determinant D is calculated as:
\vspace{-6pt}
\begin{align}
D = a_{11}
\begin{vmatrix}
a_{22} & a_{23} \\ 
a_{32} & a_{33} 
\end{vmatrix}
- a_{12}
\begin{vmatrix}
a_{21} & a_{23} \\ 
a_{31} & a_{33}  
\end{vmatrix} 
+ a_{13}
\begin{vmatrix}
a_{21} & a_{22} \\ 
a_{31} & a_{32}  
\end{vmatrix} 
\end{align}
where the $2\times2$ determinants are called \textit{minors}.
\begin{marginfigure}[10\baselineskip]
	\centering
	\includegraphics[width = 2.4in]{figures/part1b/linear_analytical/KCL.png}
	\centering
	\caption{Three loop current network}
	\label{fig:2}
\end{marginfigure}
\begin{example}
Consider a three loop current network with five resistors and two voltage sources as shown in Figure \ref{fig:2}.The electrical network consists of 5 resistors (R), two voltage sources (v) and three unknown currents in each loop (i). Given that $R_1$ = 3 $\Omega$, $R_2$ = 5 $\Omega$, $R_3$ = 2 $\Omega$, $R_4$ = 3 $\Omega$, $R_5$ = 8 $\Omega$, $v_1$ = 5 V and $v_2$ = 3 V. Set up a system of linear equations for this problem which will be solved for the three unknowns $i_1$, $i_2$ and $i_3$ using Cramer's rule. First evaluate determinant of the coefficient matrix for this system by hand and using MATLAB.
\textit{Solution}: Applying Kirchoff's current law and Ohm's law for the three current loops we get the following.
\begin{align}
-3 + 3i_1 + 5(i_1 - i_2) &= 0\\
5(i_2 - i_1) + 2i_2 + 3(i_2 - i_3)&= 0\\
3(i_3 - i_2) + 8i_3 + 3 &= 0
\end{align}
These equations can be simplified to obtain the following system of linear equations.
\begin{align}
8i_1 - 5i_2  &= 5\\
-5i_1  + 10i_2 - 3i_3 &= 0\\
-3i_2 + 11i_3 &= -3
\end{align}
\begin{enumerate}
\item Manually
\[
D = 8
\begin{vmatrix}
10 & -3 \\ 
 -3 & 11  
\end{vmatrix}
- (-5)
\begin{vmatrix}
-5 & -3 \\ 
0 & 11 
\end{vmatrix} 
+ 0
\begin{vmatrix}
-5 & 10 \\ 
0 & -3 
\end{vmatrix} 
\vspace{-2pt}
\]
\begin{align}
D = 8[(10)11 - (-3)(-3)]+5[(-5)11 - 0(-3)]+0[-5(-3) - 0(10)] = 533
\end{align}
\item MATLAB: Determinant can be computed directly in MATLAB with the \verb|det| function.
\begin{Verbatim}
>> A=[8 -5 0;-5 10 -3;0 -3 11];
>> D=det(A)
D = 533
\end{Verbatim}
\end{enumerate}
Note that the determinant of a singular system is zero and a ill-conditioned system is almost close to zero.
\textbf{Cramer's Rule}: Using this rule for a system of three linear equations, $x_1$, $x_2$ and $x_3$ would be computed as
\[
x_1 =\dfrac{
\begin{vmatrix}
b_1 & a_{12} & a_{13} \\ 
b_2 & a_{22} & a_{23} \\ 
b_3 & a_{32} & a_{33}
\end{vmatrix}
}{D},
x_2 =\dfrac{
\begin{vmatrix}
a_{11} & b_1 & a_{13} \\ 
a_{21} & b_2 & a_{23} \\ 
a_{31} & b_3 & a_{33}
\end{vmatrix}
}{D},
x_3 =\dfrac{
\begin{vmatrix}
a_{11} & a_{12} & b_1 \\ 
a_{21} & a_{22} & b_2 \\ 
a_{31} & a_{32} & b_3
\end{vmatrix}
}{D}
\]
\end{example}
\begin{example}
Use Cramer's rule to solve for the three unknowns  $i_1$, $i_2$ and $i_3$ in the above example by hand and using MATLAB.\\
\vspace{-20pt}
\begin{enumerate}
\item Manually\\
As already evaluated, determinant D = 533. Therefore
\[
i_1 =\dfrac{
	\begin{vmatrix}
	5 & -5 & 0 \\ 
	0 & 10 & -3 \\ 
	-3 & -3 & 11
	\end{vmatrix}
}{533}, i_2 =\dfrac{
	\begin{vmatrix}
	8 & 5 & 0 \\ 
	-5 & 0 & -3 \\ 
	0 & -3 & 11
	\end{vmatrix}
}{533}, i_3 =\dfrac{
	\begin{vmatrix}
	8 & -5 & 5 \\ 
	-5 & 10 & 0 \\ 
	0 & -3 & -3
	\end{vmatrix}
}{533} 
\vspace{-6pt}
\]
Which when simplified we get $i_1$ = 0.863 A $i_2$ = 0.3809 A and $i_3$ = -0.1689 A.
\item MATLAB:
\begin{Verbatim}
D=det(A);
A1 = A; A1(:,1) = b;
A2 = A; A2(:,2) = b;
A3 = A; A3(:,3) = b;
x1 = det(A1)/D;
x2 = det(A2)/D;
x3 = det(A3)/D;
\end{Verbatim}
\end{enumerate}
\begin{marginfigure}[-20\baselineskip]
	\begin{Verbatim}
A=[8 -5 0;-5 10 -3;0 -3 11];
b=[5;0;-3];
	\end{Verbatim}
\end{marginfigure}
\begin{marginfigure}[-10\baselineskip]
\textit{Output:}
\begin{Verbatim}
x1 = 0.8630
x2 = 0.3809
x3 = -0.1689
\end{Verbatim}
\end{marginfigure}
Note: For more than three equations, Cramer's rule becomes impractical because, as the number of equations increases, the determinants are time consuming to evaluate by hand (or by MATLAB).
%\end{document}
\end{example}
\section{Elimination of Unknowns}
Consider the system of two equations from graphical method as shown below
\begin{align}
3x_1 + 2x_2 &= 18\\
-x_1 + 2x_2 = 2
\end{align}
The idea behind this method is to multiply the equations by constants so that one of the unknowns will be eliminated when the two equations are combined. We will be left with one equation from which one of the unknowns is evaluated and then that is substituted back into the original equation to solve for the other unknown. Multiplying first equation with -1 and second equation with 3 we get
\begin{align}
-3x_1 - 2x_2 = -18\\
-3x_1 + 6x_2 = 6
\end{align}
Subtracting, $-8x_2 = -24$ and therefore $x_2 = 3$. Substituting this in the first equation we get,
\begin{align}
-3x_1 - 2(3) &= -18\\
x_1 &= 4
\end{align}
Note: The elimination of unknowns can be extended to systems with more than two or three equations. However, the numerous calculations that are required for larger systems make the method extremely tedious to implement by hand.

\section{Exercises}
\begin{enumerate}
    \item TODO
\end{enumerate}


\end{document}